% EGA V, English translation by P. Blass, 1987 (cote 24): archive cote 24, site page 40 (PDF page 41).
% typescript. Draft compared: none.
% AI-transcribed by Claude Opus 5.5 (claude-opus-5-5, adaptive thinking, effort high), Message Batches API, batch ega24b (msgbatch_01RGJJeHdR96U16ZWFuUjtDx), 2026-11-13; not yet checked by eye.
\scanpage{24/40}
\mspage{40}

$(T_i, X) = k$ and such that for every irreducible component $X_j$ of $X$
containing $T_i$ and such that
\[ \dim O_{X_j, T_i} = \dim O_{X, T_i} \]
i ( = k) we have
\[ X_j \over{C}{\not\subset} Y. \]

The verification of this lemma is immediate due to the general
facts in $O_{IV}$ chapter IV about dimension.  With the assumptions of
8.4, by 8.4.1(b) we see which ones are the exceptional hyperplanes
$H_\xi$.  If we exclude the set $Z_0$ of $\xi \in \underline{P}\mathrm{v}$ such that there is an
irreducible component R of T or of X such that dim f(R) > 0 and
such that $R_\xi$ is of ``dimension too large'' (a set which is of
codimension two and if what follows it does not count the
exceptional $H_\xi$ are there for which there exists a $T_i$ with codim
$(T_i, X) = k$ and dim $f(T) = 0$, $f(T)$ CH and such that for \uline{every}
irreducible component $X_j \supset T_i$ of X with codim $(T_i, X_j) = k$ we have
$f(X_j) \not\subset H_\xi$.  For a given $T_i$ with codim $(T_i, X) = k$ if there exists
an $X_j$ with codim $(T_i, X_j)$ = ? [illegible ask Grothendieck]
and such that dim $f(X_j) = 0$ then we will have
\[ f(X_j) = f(T_i)\ \mathrm{C}\ H_\xi \text{ and consequently } \xi \text{ would not be exceptional} \]
relative to the $T_i$ if and only if $f(T_i)$ C $H_\xi$ ; the set of such $\xi$ is
(the trace over) of a hyperplane of $\underline{P}\mathrm{v}$.  This proves
$\qquad\qquad P - Z_0$

8.4, and also proves the more precise result that the exceptional
set is the union of a set of codimension $\geq 2$ and of a \uline{union of
hyperplanes} determined in an evident way by the above proof.
